\begin{exercise}[Gauss law and the theta-shifted canonical momentum]{I-CH10-015}
Start with the $D=4$ Lorentzian density in mostly-plus signature,
\[
  \mathcal L=-\frac1{4g^2}F^a{}_{\mu\nu}F^{a\mu\nu}
  +\mathcal L_{\rm matter},
\]
where $\rho^a$ denotes the matter charge density with the sign convention
shown below.  Impose temporal gauge only after varying the action with respect
to $A^a{}_{0}$.  Taking $E_i{}^a$ to be the canonical momentum conjugate to
$A_i{}^a$, derive
\[
  G^a(\mathbf x)=D_i{}^{ab}E_i{}^b-\rho^a=0.
\]
For a time-independent gauge parameter that vanishes fast enough at infinity,
define $G(\omega)=\int\dd^3\mathbf x\,\omega^aG^a$.  Use canonical brackets to compute
its action on $A_i$, $E_i$, and the matter fields.  Conclude from gauge
invariance of the Hamiltonian that $[G(\omega),H]=0$, stating the boundary and
anomaly assumptions.

Finally consider an unbroken Abelian factor with a theta term normalized so
that
\[
  \Pi_i=\frac1{g^2}\mathcal E_i+\frac{\theta}{2\pi}B_i.
\]
Derive the constraint in Banks's canonical notation,
\[
  \partial_iE_i-\frac{\theta}{2\pi}\partial_iB_i-\rho=0,
\]
and explain the induced electric charge when magnetic charge is present.
\end{exercise}
